\begin{proof}
First we use the minimization analogue of FindMaxWeight to find the minimum edge weight $\alpha$.  Then by normalizing by $1/\alpha$ we may assume that all 
edge weights are at least $1$ and apply \cref{thm:rsw_alg}.  The bound on the quantum query complexities then follows from \cref{lem:compl-subroutines}.
\end{proof}



%%%%%%%%%%%%%%%%%%%%%%%%%%%%%%%%%%%%%%%%%%%%%%%%%%%%%%%%%%%%%%%%%%%
%% Time-efficient algorithm
%%%%%%%%%%%%%%%%%%%%%%%%%%%%%%%%%%%%%%%%%%%%%%%%%%%%%%%%%%%%%%%%%%%

\section{Time-efficient quantum algorithm for minimum cut} \label{sec:time-efficient}
In this section we describe a quantum algorithm for computing the weight of a minimum cut of a weighted graph with time complexity $\tO(\sqrt{mn\tau} + n^{3/2})$ in the adjacency array model and $\tO(n^{3/2}\sqrt{\tau})$ in the adjacency matrix model.
In the adjacency matrix model this is optimal up to polylogarithmic factors.
Our algorithm is a time-efficient implementation of \cref{alg:mincut}.  The running time of this algorithm is the sum of the running time of its 4 subroutines, and we have already 
analyzed the complexity of 3 of those subroutines in \cref{lem:compl-subroutines}.  Thus it now suffices to give a time-efficient implementation of the subroutine 
LearnCutAtoms, as formalized in the next lemma.

\begin{lemma}
\label{lem:kappa}
Let $\kappa(n)$ denote the maximum time complexity of a quantum algorithm for the subroutine LearnCutAtoms$(H,(1+\frac{1}{100})\lambda(H),\frac{1}{20})$ over weighted $n$-vertex graphs $H$ with $\tO(n)$ edges.
Let $G$ be a weighted graph with $n$ vertices, $m$ edges, and edge-weight ratio $\tau$.
There is a quantum algorithm to compute the weight and shores of a minimum cut of $G$ with probability at least $2/3$
that runs in time $\kappa(n) + \tO(\sqrt{mn\tau})$ in the adjacency array model and 
$\kappa(n) + \tO(n^{3/2}\sqrt{\tau})$ in the adjacency matrix model.
\end{lemma}
\begin{proof}
First we use minimum finding \cref{thm:qmin} to determine the minimum $\alpha$ and maximum $\beta$ edge weights with error probability at most $1/12$.  This requires time 
$\tO(\sqrt{m})$ in the adjacency array model and $\tO(n)$ in the adjacency matrix model and so will be low order to the time bounds stated in the lemma.  From $\alpha, \beta$ we compute the 
edge-weight ratio $\tau = \beta/\alpha$.  By multiplying all edge weights by $1/\alpha$ we may assume that the minimum edge weight is $1$ and 
the maximum edge weight is $\tau$.

If $\tau > m/n$ (in the adjacency array model) or $\tau > n$ (in the adjacency matrix model), then we simply run a randomized near-linear time algorithm (e.g., \cite{Karger00}) for calculating the weight and shores of a minimum cut of $G$.
This then takes time $\tO(m) \in \tO(\sqrt{mn\tau})$ in the array model and $\tO(n^2) \in \tO(n^{3/2}\sqrt{\tau})$ in the matrix model.
We can hence assume that $\tau \leq m/n$ in the array model and $\tau \leq n$ in the matrix model.

We use a quantum implementation of \cref{alg:mincut}.  By \cref{thm:rsw_alg} this algorithm has error probability at most $1/4$, thus our overall error probability will be at most $1/3$ as desired.  
For the running time it suffices to analyze the quantum time complexity of all 8 steps.
In \cref{lem:compl-subroutines} we show that the time complexity of steps 1--3 and 6 is $\tO(\sqrt{mn\tau})$ in the adjacency array model and $\tO(n^{3/2}\sqrt{\tau})$ in the adjacency matrix model.
For step 4, we can use a randomized near-linear time algorithm (e.g., \cite{Karger00}) for calculating the weight and shores of a minimum cut of $H$.  As $H$ has $\tO(n)$ edges this takes time $\tO(n)$.
In step 7, we compute the weight and shores of a minimum cut in $G'$ which has at most $100\tau n$ edges by the definition of LearnContraction.  This takes time $\tO(\tau n)$, which is $\tO(\sqrt{mn\tau})$ in the array model (by the assumption $\tau \leq m/n$) or $\tO(n^{3/2} \sqrt{\tau})$ in the matrix model (by the assumption $\tau \leq n$).
Finally, step 8 is trivial and the quantum time complexity of step 5 is exactly $\kappa(n)$.
\end{proof}

This section is hence devoted to proving the following theorem. 
\begin{restatable}{theorem}{cutatomsthm} \label{thm:learncutatoms}
Let $H$ be an $n$-vertex weighted graph with $m$ edges.
There is a quantum algorithm that implements LearnCutAtoms$(H,(1+\frac{1}{100})\lambda(H),\frac{1}{20})$ in time $\tO(m + n^{3/2})$.
\end{restatable}

In particular, \cref{thm:learncutatoms} implies that $\kappa(n) \in \tO(n^{3/2})$, and hence we find a time-efficient quantum algorithm.
\thmtime*

\begin{proof}
Follows from \cref{lem:kappa} and \cref{thm:learncutatoms}.  
\end{proof}

\subsection{Tools}

Our time efficient algorithm builds on a number of tools, which we first introduce here.

\subsubsection{2-respecting cuts and Karger's theorem}

In his seminal work on a near-linear time randomized algorithm for minimum cut \cite{Karger00}, Karger combined sparsification with the notion of \emph{tree-respecting cuts}.
Consider an $n$-vertex graph $G = (V,w)$, a spanning tree $T$ and a cut with shore $X$.
We say that the cut \emph{2-respects} $T$ if it cuts at most 2 edges of $T$, i.e., $|\Delta_T(X)| \leq 2$, and \emph{strictly 2-respects} $T$ if $|\Delta_T(X)| = 2$.  Note that 
the set of cuts which 2-respect $T$ depends only on $E(T)$ and not the weight of edges in $T$.  Note also that there are $n-1 + \binom{n-1}{2} = \binom{n}{2}$ cuts that 2-respect $T$.

Karger proved that we can efficiently construct a set of $O(\log n)$ spanning trees of $G$ such that every minimum cut of $G$ will 2-respect a constant fraction of them.
This effectively reduces the exponentially large search space for finding a minimum cut to the set of merely $O(n^2 \log n)$ cuts that 2-respect one of the spanning trees.
For our purpose, we will use these spanning trees as an efficient representation of the near-minimum cuts of the graph.  For this, we need a slight generalization of Karger's theorem 
on tree-respecting cuts.  This shows we can efficiently find $O(\log n)$ spanning trees such that any $(1+1/16)$-near-minimum cut $2$-respects a constant fraction of them, while Karger's statement was 
only for minimum cuts.  This only requires a minor modification of Karger's proof, but for completeness we provide a proof in \cref{app:karger}. 

Throughout this section we will use the phrase ``with high probability'' to mean with probability at least $1-1/n^c$ for an arbitrary constant $c$.
\begin{restatable}[{\cite[Theorem 4.1]{Karger00}}]{theorem}{kargerthm} \label{thm:karger}
Let $G = (V,w)$ be a weighted graph with $n$ vertices and $m$ edges.  There is a randomized algorithm that in time 
$O(m \log^2(n) + n\log^4(n))$ time constructs a set of $O(\log n)$ spanning trees such that every $(1+1/16)$-near minimum cut of $G$ 2-respects $1/4$ of them with high probability.
\end{restatable}
Karger states the runtime of the algorithm in this theorem as $O(m + n \log^3(n))$, but we opt for a simpler proof rather optimizing log factors.

\subsubsection{Data structures}
We will frequently need to refer to a 2-respecting cut both by its shores and the edges of the tree it cuts.  We develop some notation to make this easier.
\begin{definition}[Notation for 2-respecting cuts]
Let $T$ be a tree on vertex set $V$ with root $r$.  Define $N(T) = E(T) \cup E(T)^{(2)}$.
For $f \in N(T)$ define $\shore(f)$ to be the set $X \subseteq V$  such that $\Delta_T(X) = f$ and $X$ does not contain $r$.
For $X \subseteq V$ such that $|\Delta_T(X)| \le 2$, 
let $\edges(X) = \Delta_T(X)$.  We overload both these notations to sets so that $\shore(Q) = \{\shore(f): f \in Q\}$ for $Q \subseteq N(T)$ and similarly 
$\edges(\Tcal) = \{\Delta_T(X): X \in \Tcal\}$ for a set $\Tcal$ of shores of 2-respecting cuts of $T$.  
\end{definition}

With some preprocessing time, we can efficiently evaluate the weight of 2-respecting cuts.  The following lemma is very useful.
\begin{lemma}[{\cite[Lemma 1]{GMW20c}}] \label{lem:eval-cuts}
Given a weighted graph $G = (V,w)$ with $n$ vertices and $m$ edges, and a spanning tree $T$ of $G$, we can construct in $O(m \log n)$ time a data structure that, for any $f \in N(T)$,
reports the weight $w(\Delta_G(\shore(f)))$ of the corresponding 2-respecting cut in $O(\log n)$ time.
\end{lemma}



Another data structure that we use is based on the \emph{Euler tour technique} \cite{TV84,HK95}.
This is a way of representing a tree that is useful to access and modify data in subtrees.
Consider an undirected tree $T=(V_T,E_T)$ with root $r \in V_T$.
To $T$ we associate the directed graph $\vec T = (V_T, \vec E_T)$ obtained by replacing every edge in $E_T$ by a pair of directed edges in opposite directions.
Now let $\Ecal_T \in (\vec E_T)^{2(n-1)}$ denote an Euler tour in $\vec T$, starting and ending in root $r$.
$\Ecal_T$ is a sequence of $2(n-1)$ edges as each directed edge is traversed exactly once.

For every node $u$ in $V_T$, let $f(u)$ be the index in $\Ecal_T$ of the edge that points toward $u$, and let $\ell(u)$ be the index of the last edge that points toward $u$.
Now if $T(u)$ is the subtree of $T$ induced by vertex $u$ and all of its descendants, then the subsequence of $\Ecal_T$ starting at $f(u)$ and ending at $\ell(u)$ (both included) is an Euler tour representation of $T(u)$.
Hence any subtree corresponds to a subsequence of $\Ecal_T$.
We can use this to prove the lemma below, which will be useful to compute $\atoms(\Tcal)$ from a set $\Tcal$ of shores of cuts that 2-respect a given tree.

Given a tree whose nodes have some key value, we call a \emph{subtree-add} the increasing or decreasing of the key value in a subtree by some fixed amount.

\begin{lemma} \label{lem:euler-tree}
Let $T=(V_T,E_T)$ be a tree with key values $\{k_u \mid u \in V_T\}$ of $O(\log n)$ bits.
There is a data structure that implements $M$ subtree-adds in time $\tO(n + M)$.
\end{lemma}
\begin{proof}
Fix a root node $r$.
Represent $T$ by an Euler tour $\Ecal_T \in (\vec E_T)^{2(n-1)}$ and define $f(u),\ell(u)$ for each $u \in V$ as above.
Associate to $\Ecal_T$ a list $A$ of length $2(n-1)$ to store the key values, setting $A(i) = k_u$ if the $i$-th entry of $\Ecal_T$ is an edge whose tail is $u$.
Adding value $\alpha$ to the keys of nodes in subtree $T(u)$ amounts to adding $\alpha$ to every entry in the subsequence in $A$ starting with $f(u)$ and ending with $\ell(u)$ (both included).
Call such an operation $\texttt{ADD}(\alpha,f(u),\ell(u))$.

To implement $M$ \texttt{ADD} operations, create a second emtpy list $B$ with length $2(n-1)$.
For every operation $\texttt{ADD}(\alpha,f(u),\ell(u))$, set $B(f(u)) = B(f(u)) + \alpha$ and if $\ell(u) < 2(n-1)$ set $B(\ell(u)+1) = B(\ell(u)+1) - \alpha$.
Now do a partial sum transformation of $B$:
\begin{algorithmic}[1]
\State
Create list $s_B$ of length $2(n-1)$ with $s_B(1) = B(1)$ and $s_B(i) = 0$ for all $i \in [2,2(n-1)]$.
\For{$i = 2,3,\dots,2(n-1)$}
\State
Set $s_B(i) = s_B(i-1) + B(i)$.
\EndFor
\end{algorithmic}
In total this has time complexity $\tO(n + M)$ (assuming $\tO(1)$ cost for arithmetic operations).
The final key values are now given by setting $k_u = A(f(u)) + B(f(u))$.
\end{proof}

\subsection{Generating set for a single tree} \label{sec:single-tree}
Let $G = (V,w)$ be an $n$-vertex weighted graph and $T$ be a spanning tree of $G$.  Let $Q \subseteq E(T)^{(2)}$ and 
$\Mcal = \shore(Q)$.  In words, $\Mcal$ is an arbitrary set of shores of cuts that \emph{strictly} 2-respect $T$.  
The next lemma gives an explicit generating set $\Scal$ for $\atoms(\Mcal)$ with $|\Scal| \le n-2$.  We first make a definition
that will be used throughout this section.

\begin{definition}[separate]
Let $V$ be a finite set and $X \subseteq V$.  For $u,v \in V$ we say that $X$ separates $u,v$ if exactly one of 
them is in $X$.
\end{definition}

\begin{lemma}
\label{lem:spanning}
Let $T$ be a tree on a vertex set $V$ of cardinality $n$.  Let $\Mcal \subseteq 2^V$ be a set of shores that strictly 2-respect $T$ and let 
$Q = \edges(\Mcal)$.  Define the graph $L = (E(T), Q)$ and let $F$ be a spanning forest of $L$.  
Then $\Scal = \shore(E(F))$ is a generating set for $\atoms(\Mcal)$.
\end{lemma}

\begin{proof}
Clearly $E(F) \subseteq Q$ thus $\Scal \subseteq \Mcal$.  This means that $\atoms(\Mcal)$ is a refinement of $\atoms(\Scal)$.  
Thus to show $\atoms(\Scal) = \atoms(\Mcal)$ it suffices to show that any $u,v \in V$ that are in different sets of $\atoms(\Mcal)$ are also in 
different sets of $\atoms(\Scal)$.  
 
The key fact we need is that if $\Delta_T(X) = \{e,e'\}$ then $X$ separates $u,v$ iff exactly one of $e,e'$ is on the 
path from $u$ to $v$ in $T$.  Suppose that $u,v$ are in different sets of $\atoms(\Mcal)$, that is there is an $X \in \Mcal$ 
which separates them.  Say that $\Delta_T(X) = \{e_\rmin,e_\rmout\}$ where $e_\rmin$ is on the $u-v$ path in 
$T$ and $e_\rmout$ is not.  Then $\{e_\rmin,e_\rmout\} \in Q$ and therefore there must be a path between $e_\rmin$ and $e_\rmout$ in the spanning forest 
$F$.  Let $(e_0,e_1, e_2, \ldots, e_k)$, where $e_0 = e_\rmin, e_k = e_\rmout$, be the sequence of vertices on this path in $F$.  As $e_\rmin$ is on the $u-v$ path in $T$ and $e_\rmout$ is 
not, there must be consecutive vertices $e_i,e_{i+1}$ where $e_i$ is on the $u-v$ path in $T$ and $e_{i+1}$ is not.  As $\{e_i, e_{i+1}\} \in E(F)$ there is 
an $X \in \Scal$ which separates $u$ and $v$.  
\end{proof}

\begin{lemma}
\label{lem:q_genset}
Let $G = (V,w)$ be an $n$-vertex weighted graph with $m$ edges and $T$ a spanning tree of $G$.  For a real number 
$\alpha \ge 1$, let $\Tcal = \{X \subseteq V: w(\Delta_G(X)) \le \alpha \lambda(G), |\Delta_T(X)| \le 2\}$.  
There is a quantum algorithm that outputs with high probability a set $Q \subseteq N(T)$ in time $\tO(m + n^{3/2})$ such that $|Q| \le 2n-3$ and $\Scal = \shore(Q)$
is a generating set for $\atoms(\Tcal)$.  
\end{lemma}

\begin{proof}
Let $\Tcal_1 = \{X \in \Tcal: |\Delta_T(X)| =1\}$ and $\Tcal_2 = \{X \in \Tcal: |\Delta_T(X)| =2\}$. Let $Q_1 = \edges(\Tcal_1)$ and $Q_2 = \edges(\Tcal_2)$.  
Let $F$ be a spanning tree for $L = (E(T), Q_2)$.  By \cref{lem:spanning}, $\Rcal = \shore(E(F))$ is a generating set for 
$\atoms(\Tcal_2)$ and $|\Rcal| \le n-2$ as $F$ is a spanning tree of an $n-1$-vertex graph.  Thus by \cref{prop:atom_cup}, $\Scal = \Tcal_1 \cup \Rcal$ is a generating set for $\Tcal$ of size at most $2n-3$.  
Thus taking $Q = Q_1 \cup E(F)$ satisfies the conditions of the lemma.

Now we must show how to efficiently output $Q$.  We can first run a near-linear time classical randomized algorithm to compute $\lambda(G)$ \cite{Karger00}.  We then in near-linear time set up the 
data structure given by \cref{lem:eval-cuts}.  For an $f \in N(T)$ this lets us check in time $O(\log(n))$ if $f \in Q_1 \cup Q_2$.  We can then cycle over the edges $e\in E(T)$ to create the set $Q_1$ classically in time 
$\tO(n)$.  It now remains to construct a spanning tree of $L = (E(T), Q_2)$.
For any $f \in E(T)^{(2)}$ we can use the data structure to check in $O(\log n)$ time if $f \in Q_2$.
This gives us adjacency matrix access to $L$ with $O(\log n)$ overhead for each query.  
Now we can use the quantum algorithm from \cite{DHHM06} that with high probability outputs a spanning forest of an $n$ vertex graph in the adjacency matrix model with $\tO(n^{3/2})$  queries and time.  
Thus we can use this algorithm to construct a spanning forest $F$ of $L$.  We then output $Q = Q_1 \cup E(F)$ as desired.
\end{proof}

Now we have an implicit representation $\edges(\Scal)$ of a generating set $\Scal$ for $\atoms(\Tcal)$, where $\Tcal$ is the set of near-minimum cuts of a graph $G$ that 2-respect a tree $T$.  
What we need, however, is to actually output $\atoms(\Tcal)$.
In the following lemma we show how to do this efficiently by combining random hashing with Euler tour trees.

\begin{lemma} \label{lem:2-resp-part}
Let $T$ be a tree on a vertex set $V$ of size $n$, $Q \subseteq N(T)$, and $\Scal = \shore(Q)$.  
Given input $Q$ there is a classical algorithm that with probability at least $1-1/n$ outputs $\atoms(\Scal)$ in time $\tO(n + |Q|)$.
\end{lemma}
\begin{proof}
Let $M$ be a large integer to be chosen later and consider the following algorithm.
Pick $\ell \in \Z_M$ uniformly at random and give every vertex $u \in V$ the key value $k_u = \ell$.
For every $f \in Q$, do:
\begin{itemize}
\item
Pick $\ell \in \Z_M$ uniformly at random and set $k_u = k_u + \ell \;(\bmod M)$ for all $u \in \shore(f)$.
\end{itemize}
Now if $u$ and $v$ are in the same set of $\atoms(\Scal)$, that is no set of $\Scal$ separates them, then $k_u = k_v$.
On the other hand, if $u$ and $v$ are in different sets of $\atoms(\Scal)$ then there is some $f \in Q$ such that $u \in \shore(f)$ and $v \not \in \shore(f)$, or vice versa. In this case,
$k_u$ and $k_v$ are pairwise independent 
and distributed uniformly at random in $\Z_M$.
Hence $k_u = k_v$ with probability $1/M$.
Taking a union bound over all pairs $u,v$, we see that with probability at least $1-\binom{n}{2}/M$ we have that $k_u \neq k_v$ for all $u,v$ in different sets of $\atoms(\Scal)$.
If we set $M = n^3$ and we let $\Pcal(\{k_u\})$ denote the partition induced by gathering nodes with the same key value, then $\Pcal(\{k_u\}) = \atoms(\Scal)$ with probability at least $1 - 1/n$.

The cost of actually implementing this algorithm is dominated by sequentially updating for every $f \in Q$ the key value for all nodes in $\shore(f)$.
This amounts to changing the key value in at most 2 subtrees of $T$:
\begin{itemize}
\item
If $f = e \in E(T)$, then $\shore(f)$ is the subtree $T(u)$ of some node $u$ and we have to change the key value in $T(u)$.
\item
If $f = \{e,e'\} \in E(T)^{(2)}$, then we distinguish two cases.
If one of the two cut edges is a descendant of the other then $\shore(f)$ is of the form $T(u) \setminus T(v)$ for two nodes $u,v \in V$.
In this case we can update the key values by adding $\ell$ to $T(u)$ and subtracting $\ell$ from $T(v)$.
If neither of the edge is a descendant of the other then $\shore(f)$ is of the form $T(u) \cup T(v)$, and we can update the key values by adding $\ell$ to $T(u)$ and $T(v)$.
\end{itemize}
In \cref{lem:euler-tree} we show how to change the key values in $|Q|$ subtrees in total time $\tO(n + |Q|)$ using Euler tour trees.
\end{proof}

We can now put all these pieces together into the following algorithm.
\begin{algorithm}[H]
\caption{Algorithm for finding atoms of the shores of 2-respecting near-minimum cuts}
\label{alg:partition-tree}
 \hspace*{\algorithmicindent} \textbf{Input:} Explicit description of $G=(V,w)$, a spanning tree $T$ of $G$, a real number $\alpha \ge 1$. \\
 \hspace*{\algorithmicindent} \textbf{Output:} $\atoms(\Tcal)$ where $\Tcal = \{X: w(\Delta_G(X)) \le \alpha \lambda(G) \mbox{ and } |\Delta_T(X)| \le 2\}$.
\begin{algorithmic}[1]
\State
Compute $\lambda(G)$.
\State
Create data structure as in \cref{lem:eval-cuts} for evaluating the weight of cuts in $G$ that 2-respect $T$.
\State
Compute $Q$ such that $\shore(Q)$ is a generating set for $\atoms(\Tcal)$ by \cref{lem:q_genset}.
\State
Use \cref{lem:2-resp-part} to find and return $\atoms(\shore(Q)) = \atoms(\Tcal)$.
\end{algorithmic}
\end{algorithm}

\begin{lemma} \label{lemma:2-respect}
Let $G=(V,w)$ be an $n$-vertex weighted graph with $m$ edges and $T$ a spanning tree of $G$.  Let $\alpha \ge 1$ be a real number and
$\Tcal = \{X: w(\Delta_G(X)) \le \alpha \lambda(G) \mbox{ and } |\Delta_T(X)| \le 2\}$.  
\cref{alg:partition-tree} outputs $\atoms(\Tcal)$ with high probability and can be implemented by a quantum algorithm in time $\tO(m + n^{3/2})$.
\end{lemma}

\subsection{Time-Efficient quantum algorithm for LearnCutAtoms}
We now describe a time-efficient quantum algorithm for outputting $\atoms(\Tcal)$, where $\Tcal$ is the set of shores of all $(1+1/100)$-near-minimum cuts of a weighted graph $H$.
This algorithm combines Karger's tree packing \cref{thm:karger} with the algorithm that produces the atoms of shores of cuts that 2-respect a tree from the previous section (\cref{lemma:2-respect}).

\begin{algorithm}[H]
\caption{LearnCutAtoms$(H,\lambda,\delta)$}
\label{alg:partition}
 \hspace*{\algorithmicindent} \textbf{Input:} Explicit description of an $n$-vertex weighted graph $H = (V,w)$ with $m$ edges, a cut threshold $\lambda \le (1+1/16) \lambda(H)$, and an error parameter $\delta$. \\
 \hspace*{\algorithmicindent} \textbf{Output:} $\atoms(\Tcal)$ where $\Tcal = \{X \subseteq V: w(\Delta_G(X)) \le \lambda\}$.
\begin{algorithmic}[1]
\State
Construct set of $K \in O(\log n)$ spanning trees $\{T_i\}$ using \cref{thm:karger}.
\For{$i = 1,2,\dots,K$}
\State
Use \cref{alg:partition-tree} to find $\atoms(\Tcal_i)$ where $\Tcal_i = \{X \subseteq V: w(\Delta_G(X)) \le \lambda \mbox{ and } |\Delta_{T_i}(X)| \le 2\}$.
\EndFor
\State
Output $\atoms(\cup_i \atoms(\Tcal_i))$.
\end{algorithmic}
\end{algorithm}

\cutatomsthm*

\begin{proof}
We use \cref{alg:partition}.  First let us argue correctness.  As $\lambda \le (1+16)\lambda(H)$, by \cref{thm:karger} for every $X \in \Tcal$ 
there will be a tree $T_i$ such that $\Delta_{T_i}(X) \le 2$.  This means that $\Tcal = \cup_{i=1}^K \Tcal_i$.  Hence 
$\atoms(\Tcal) = \atoms(\cup \Tcal_i) = \atoms(\cup \atoms(\Tcal_i))$.  By \cref{lemma:2-respect}, step~(3) correctly outputs $\atoms(\Tcal_i)$ 
for $i = 1, \ldots, K$ with high probability, and thus step~(5) will output $\atoms(\Tcal)$ with high probability.

Now let us analyze the complexity.  Step~(1) can be done in $\tO(m)$ time by a classical randomized algorithm by \cref{thm:karger}.  Step~(3) can 
be done by a quantum algorithm in time $\tO(m + n^{3/2})$ by \cref{lemma:2-respect}, and thus the for loop has the same time bound as $K = O(\log n)$.

Finally, we need to explain how to (classically) implement step~(5).  First we give every node $v \in V$ a key value $k_v = 0$.
Then, for each $i=1,\ldots, K$, we iterate over the node set and append a $\log n$-bit string to the key value of every node, indicating the component of $\atoms(\Tcal_i)$ of which it is part.
At the end of this routine every node has a $O(\log^2 n)$-bit key value that indicates its component in $\atoms(\Tcal)$.  The total runtime for this step is $\tO(n)$.  Thus overall 
the running time is $\tO(m + n^{3/2})$.
\end{proof}



%%%%%%%%%%%%%%%%%%%%%%%%%%%%%%%%%%%%%%%%%%%%%%%%%%%%%%%%%%%%%%%%%%%
%% Lower bounds
%%%%%%%%%%%%%%%%%%%%%%%%%%%%%%%%%%%%%%%%%%%%%%%%%%%%%%%%%%%%%%%%%%%

